\documentclass[12pt,oneside]{book}
\usepackage{geometry}
 \geometry{
 left=3.5cm,
 right=2.5cm,
 top=3cm,
 bottom=2.8cm,
 headheight=1cm,
 headsep=0.5cm,
 footskip=0.6cm,
 }
\linespread{1.5}
\usepackage{fancyhdr}
\fancypagestyle{plain}[fancy]{}
\renewcommand{\headrulewidth}{0pt}
\usepackage{lipsum} 

\usepackage{titlesec}
\titleformat{\chapter}[block]
  {\normalfont\huge\bfseries}{\thechapter .}{20pt}{\huge}
\usepackage{xcolor}
\usepackage[hidelinks]{hyperref}
\hypersetup{
colorlinks,
    linkcolor={blue!80!black},
    citecolor={blue!80!black}
}
\usepackage{setspace}

\usepackage{amsmath,amssymb,amsthm}
\usepackage{mathtools}

\theoremstyle{definition}
\newtheorem{definition}{Definition}[section]
\newtheorem{theorem}{Theorem}[section]
\newtheorem{proposition}{Proposition}[section]
\newtheorem{lemma}[theorem]{Lemma}
\newcommand{\bC}{\mathbb{C}}
\newcommand{\RR}{\mathbb{R}}
\newcommand{\bN}{\mathbb{N}}
\newcommand{\ZZ}{\mathbb{Z}}
\DeclareMathOperator{\GL}{GL}


\begin{document}

\thispagestyle{empty}


%PORTADA: Por favor completar la siguiente información, respetando el uso de mayúsculas tal como se indica.

\pagestyle{fancy}
\fancyfoot[c]{\thepage}
\frontmatter

\begin{center}


{\Large UNIVERSIDAD DE COSTA RICA}

{\Large SISTEMA DE ESTUDIOS DE POSGRADO}

\vspace{3cm}
{\Large E-POLYNOMIALS OF CHARACTER VARIETIES OF FINITELY GENERATED GROUPS} %Todo en mayúscula

\vspace{3cm}
{\Large Thesis submitted to the consideration of the Committee of the Graduate Studies Program in Mathematics to obtain the degree and title of Academic Master's in Mathematics with an emphasis in Pure Mathematics  } 

\vspace{3cm}
{\Large ALLAN FRANCISCO MURILLO BADILLA} %Todo en mayúscula


\vspace{3cm}
{\large Ciudad Universitaria Rodrigo Facio, Costa Rica}

\vspace{1.5cm}
{\large 2026}


\end{center}


\newpage


\chapter*{Dedication}
\addcontentsline{toc}{section}{Dedicatoria}

Digite aquí la dedicatoria, si así lo desea; de lo contrario, puede eliminar esta página.  Puede unificar la dedicatoria y los agradecimientos en una misma sección.


\newpage


\chapter*{Acknowlegments}
\addcontentsline{toc}{section}{Agradecimientos}

Digite aquí los agradecimientos, si así lo desea; de lo contrario, puede eliminar esta página.  Puede unificar la dedicatoria y los agradecimientos en una misma sección.


\newpage
% No se debe colocar ningún título, como “Hoja de aprobación” o similares (en el caso de las defensas virtuales ver las directrices del SEP). En el nombre de la persona sustentante no debe figurar ninguna indicación al grado académico que ya posea; esta página debe contener de manera exacta solamente la siguiente información:

\refstepcounter{section}
\addcontentsline{toc}{section}{Approval Page}

\begin{center}
``This thesis was accepted by the Mathematics Graduate Program Committee of the University of Costa Rica, in partial fulfillment of the requirements for the degree of Professional Master in Mathematical Methods and Applications.''

\begin{spacing}{1}
\vspace{1.5cm}
\underline{\hspace{10cm}} \\
Grado académico y nombre completo \\
\textbf{Decano(a) o Representante del/de la Decano(a)
Sistema de Estudios de Posgrado}


\vspace{1.5cm}
\underline{\hspace{10cm}} \\
Grado académico y nombre completo \\
\textbf{Profesor(a) Guía}



\vspace{1.5cm}
\underline{\hspace{10cm}} \\
Grado académico y nombre completo \\
\textbf{Lector(a)}


\vspace{1.5cm}
\underline{\hspace{10cm}} \\
Grado académico y nombre completo \\
\textbf{Lector(a)}


\vspace{1.5cm}
\underline{\hspace{10cm}} \\
Grado académico y nombre completo \\
\textbf{Director(a)/Representante Programa de Posgrado en Matemática}


\vspace{1.5cm}
\underline{\hspace{10cm}} \\
Nombre completo \\
\textbf{Sustentante}
% NO debe hacer indicación del grado académico de la persona sustentante.
\end{spacing}

% Importante: deberá utilizar el género que corresponda en las líneas de los miembros(as) del tribunal examinador, deberá utilizar -por ejemplo- “Asesor” o “Asesora” según sea el caso. Asimismo deberá indicar solamente un cargo, por ejemplo “Decana” o “Representante de la Decana”. La cantidad de miembros del tribunal examinador a incluir en la hoja de aprobación, se determina según el acta correspondiente a la presentación oral y pública (examen de grado) del estudiante.

\end{center}

\newpage
\thispagestyle{fancy}
\fancyhead{\empty}
\tableofcontents



\newpage
\chapter*{Resumen}
\addcontentsline{toc}{section}{Resumen}

Agregar en esta sección el \textbf{resumen en español}.  Este resumen es necesario independientemente del idioma del trabajo.  

La extensión del resumen debe ser de una página como máximo. Es la única que puede presentarse con interlineado sencillo de ser necesario. Debe constar únicamente, el título “Resumen” y la síntesis del documento.





\newpage
\chapter*{Abstract (o resumen en otro idioma)}
\addcontentsline{toc}{section}{Abstract}

Agregar en esta sección el abstract o resumen en otro idioma en caso de que el cuerpo del trabajo sea en un idioma distinto al español.  Debe cumplir con las características de extensión y formato del resumen en español.




\newpage
\addcontentsline{toc}{section}{Lista de tablas}
\listoftables
Agregar acá la lista de tablas cuando sea necesario.  De lo contrario, puede eliminar esta página.


\newpage
\addcontentsline{toc}{section}{Lista de figuras}
\listoffigures
Agregar acá la lista de figuras cuando sea necesario.  De lo contrario, puede eliminar esta página.


\newpage
\addcontentsline{toc}{section}{Abreviaturas}
\chapter*{Abreviaturas}
Agregar acá las abreviaturas cuando sea necesario.  De lo contrario, puede eliminar esta página.



\mainmatter

\pagestyle{fancy}
\fancyhead[r]{\thepage}
\fancyfoot[c]{\empty}





\newpage
\chapter{Introducción}
\section{Justificación}
\section{Objetivos}
\section{Otros}


\newpage
\chapter{Preliminaries}

\section{Conventions and notation}

Tangent space of a smooth manifold
Vector fields




\section{Lie Groups and its Lie Algebras}

A Lie group $G$ is a group with the additional structure of a smooth manifold, such that the group operations of multiplication and inversion are smooth maps. Lie groups always admit an analytic atlas, unique up to analytic diffeomorphism, such that multiplication and inversion are analytic maps (c.f.\ \cite[Prop 1.117]{KnappAnthonyW.1996LGBa}). A homomorphism $\varphi : G_1 \to G_2$ of Lie groups is a map that is both compatible with the group structures and smooth as a map of manifolds.


\begin{definition}
	A Lie subgroup of a Lie group G is an injective homomorphism of Lie groups i: H--, G.
\end{definition}


The standard examples of Lie groups are $\GL_n(\RR)$, $\GL_n(\bC)$, and all their closed subgroups, which are called linear Lie groups.


Let $g$ be an arbitrary element of the Lie group $G$. The symbol $L_g$ denotes left translation in $G$ by $g$, that is, $L_g(h)=gh$ for all $y \in G$. The defining properties of $G$ imply that $L_g$ is a diffeomorphism, with inverse $(L_g)^{-1}$, and the derivative $d(Lx) : TG \to TG$ must also be a diffeomorphism. For each tangent space $T_hG$, ${d L_g: T_h G \to T_h G} $ is a linear isomorphism.

\begin{definition}
    The vector field $X$ on $G$ is said to be left-invariant if and only if $d(L_g)X= X \circ L_g$ , for all points $g \in G$,
\end{definition} 
Such a field $X$ is completely determined by its value at the identity $X_e$. Conversely for every vector $v \in T_e(G)$ we can define a (smooth) left-invariant vector field $X$ on $G$ by $X_g = dL_g(v)$. 

\begin{proposition}
The map $X  \to X_e$ is a linear isomorphism between the set of left-invariant vector fields on $G$ and $T_eG $, the tangent space at the identity.
\end{proposition}


The Lie algebra of the group $G$, denoted by $\mathfrak{g}$, is the space of all left-invariant vector fields on $G$ endowed with the Lie bracket arising from the commutator (see \cite[Chapter 2, Section 3]{milneLAG}). As we saw before, we can identity the Lie Algebra $\mathfrak{g}$ with the tangent space at the identity. 

Now we recall the definition of a one-parameter group:

\begin{definition}
    A one-parameter group of diffeomorphisms of a manifold $M$ is a smooth map $\rho : M \times \RR \to M$ such that, for all $x \in M$ and $s, t \in \RR$,
    $\rho(x, s + t) = \rho(\rho(x, t), s).$
\end{definition}

A one-parameter group can be thought of as a family of curves for each point $x \in M$, so by taking derivatives at each point we can induce a global vector field. The converse is locally true by the fundamental theorem of ordinary differential equations: every smooth vector field on $M$ generates a local one-parameter group of diffeomorphisms (its flow). If the vector field is complete, this one-parameter group is defined globally over all $\RR$.

In the context of a Lie group $G$, it is a fundamental fact that every left-invariant vector field is complete (see \cite{. Therefore, every element $X \in \mathfrak{g}$ generates a unique global one-parameter subgroup of $G$, which is a smooth homomorphism $\gamma: \RR \to G$ such that $\gamma'(0) = X$. This correspondence naturally leads to the definition of the exponential map.

Definition 2.7.4. Defining Pt(x) = p(x, t), we note that this means that Ps o Pt= Ps+t and
this immediately implies that Po is the identity and then that Pt - I = P-t·
Thus we do indeed have a group of diffeomorphisms of M, isomorphic with the
additive group of real numbers.
Conversely, given such a I-parameter group of diffeomorphisms, we may
re-interpret it as a family of curves 'Yx, defined for each x E U by 'Yx(t) = Pt(x),
and then recover the vector field of which these are the integral curves as
X(x) = 'Yx'(O). For then, if y = 'Yx(t), the group property becomes ')'y = 'Yx o Tt
and then we have X(y) = 'Y/(0) = 'Yx'(t) as required for 'Yx to be an integral
curve.
Although a I-parameter group of diffeomorphisms always determines a
vector field in this way, which is said to generate the group, vector fields in
general only determine a local version of a I-parameter group. A suflicent
condition for a full group is as follows.
To connect the Lie algebra $\mathfrak{g}$ back to the Lie group $G$, we introduce the exponential map. For every left-invariant vector field $X \in \mathfrak{g}$, there exists a unique one-parameter subgroup, which is a smooth homomorphism $\gamma_X : \RR \to G$, such that its initial tangent vector is $X_e$. The Lie exponential map $\exp : \mathfrak{g} \to G$ is defined by evaluating this subgroup at $t = 1$, that is,
\begin{equation}
    \exp(X) = \gamma_X(1).
\end{equation}
This map is a local diffeomorphism between a neighborhood of the origin in $\mathfrak{g}$ and a neighborhood of the identity $e$ in $G$. 

In the specific case of linear Lie groups, the Lie exponential map coincides with the standard matrix exponential. For any matrix $X$ in the corresponding Lie algebra, it is given by the absolutely convergent series:
\begin{equation}
    \exp(X) = e^X = \sum_{k=0}^{\infty} \frac{X^k}{k!}.
\end{equation}

The group $G$ acts on itself by conjugation, $C_g(h) = g h g^{-1}$. By differentiating this action at the identity $e \in G$, we obtain the adjoint representation of the Lie group on its Lie algebra, denoted by $\operatorname{Ad}: G \to \operatorname{Aut}(\mathfrak{g})$. Explicitly, for $g \in G$ and $\xi \in \mathfrak{g}$, it is given by
\begin{equation}
    \operatorname{Ad}(g)(\xi) = \frac{d}{dt}\bigg|_{t=0} g \exp(t\xi) g^{-1}.
\end{equation}

We can classify a Lie algebra $\mathfrak{g}$ under the following categories \cite{maret_note_nodate}:
\begin{itemize}
    \item A Lie algebra $\mathfrak{g}$ is \textbf{simple} if it is not abelian and if its only proper ideal is the zero ideal. Equivalently, $\mathfrak{g}$ is simple if and only if its adjoint representation is irreducible and $\mathfrak{g}$ is not a one-dimensional abelian Lie algebra.
    \item A \textbf{semisimple} Lie algebra $\mathfrak{g}$ is one without nonzero abelian ideals; equivalently, it is a direct sum of simple Lie algebras. The Killing form $K(v_1, v_2) := \operatorname{Tr}(\operatorname{ad}(v_1) \operatorname{ad}(v_2))$ for $v_1, v_2 \in \mathfrak{g}$, where $\operatorname{Tr}$ is the endomorphism trace, provides a symmetric bilinear form that is non-degenerate precisely when $\mathfrak{g}$ is semisimple.
    \item A \textbf{reductive} Lie algebra is a direct sum of a semisimple and an abelian algebra, or equivalently, one whose adjoint representation is completely reducible.
\end{itemize}

We call a connected Lie group simple, semisimple or reductive if its Lie algebra is simple, semisimple or reductive, respectively. Simple Lie groups are semisimple and semisimple Lie groups are reductive.

\section{Algebraic Groups}



\section{}


Definitions 8.3.3. For a group G acting on a set X, the isotropy subgroup
at x EX is Gx = {g E Glgx = x}. The orbitofx is Gx = {gx I g E G}.
The action is said to be transitive if it has exactly one orbit. When X is
a topological space, the orbit space is the quotient space: the set of orbits
topologised with the quotient topology.
We note that, with each orbit, we can associate a conjugacy class of
isotropy subgroups, since the isotropy subgroup at g(x) is gGxg- 1 , and that,
although our notation is for left actions, we can equally well speak of right
actions of G on X. A left action gives rise to a right action via the rule
xg = g- 1 x, and conversely.
The quotient or homogeneous space G / H, comprising cosets of the form
xH, can be thought of as the space of orbits for the right action of H on G
given by (g, h) t---. gh. The group G itself acts on G/ H according to the rule,
the left action,
G x G/H--+ G/H;
(g,xH) t---. gxH.
We wish to show that, subject to the choices we have made, and reverting to
the language of Chapter 3, G has the structure of a principal H-bundle over
G/ H. Note, before we start the main proof, that H operates on the total space
G on the right, but that in changing charts we multiply by 9ij(x) EH on the
left. These two actions commute.
Theorem 8.3.4. Let G be a Lie group and H a closed subgroup of G. Then
G admits the structure of an H -principal bundle with projection p : G ---.
G / H; gt---. gH, p is a submersion of rank equal to the dimension of G / H, and
G / H inherits the structure of a differential manifold.

\begin{definition}
    An affine algebraic group $G$ is an algebraic variety (the common zero set of finitely many polynomials) equipped with a group structure such that multiplication and inversion, $m : G \times G \to G$, $(g, h) \mapsto gh$, and $i : G \to G$, $g \mapsto g^{-1}$, are regular morphisms of varieties.
\end{definition}

Every real or complex algebraic group also has a natural Lie group structure \cite[Chapter 3, Theorem 2.1]{milneLAG}.

\section{Finitely Generated Groups}

\begin{definition}
    A group $\Gamma$ is finitely generated if it admits a finite set of generators $S$ such that every element of $\Gamma$ is a finite product of elements of $S$ and their inverses.
\end{definition}

One of the main examples of finitely generated groups are the surface groups. 

\begin{definition}[Surface Groups]
    Let $g, n \geqslant 0$. The surface group of genus $g$ with $n$ punctures is 
    \begin{equation}
        \pi_{g,n} := \left\langle a_1, b_1, \dots, a_g, b_g, c_1, \dots, c_n \;\middle|\; \prod_{i=1}^g [a_i, b_i] = \prod_{j=1}^n c_j \right\rangle,
    \end{equation}
    where $[a_i, b_i] = a_i b_i a_i^{-1} b_i^{-1}$ for $i = 1, \dots, g$. 
    
    When $n = 0$, this corresponds to the fundamental group of a closed surface of genus $g$. For $n \geqslant 1$, the relation implies $\pi_{g,n}$ is a free group of rank $2g + n - 1$. The groups $\pi_{0,0}$ and $\pi_{0,1}$ are trivial. Only $\pi_{0,2} \cong \ZZ$ and the torus group $\pi_{1,0} \cong \ZZ^2$ are abelian; higher genus groups are non-abelian.
\end{definition}

The following theorem justifies the name \emp{surface group}

\begin{theorem}
    Let $\Sigma_{g,n}$ be an orientable surface of genus $g$ with $n$ punctures. Then $\pi_1(\Sigma_{g,n}) \cong \pi_{g,n}$.
\end{theorem}

The proof can be checked on \cite{casimiroCharacterVarietiesTheir2026} Theorem 1.


\newpage
\chapter{Character varieties}
\section{Representation variety}



\section{GIT quotient}



\section{Character variety}

\newpage

\chapter{Hodge structures}
\input{hodge_structures.tex}

\newpage
\chapter{E-polynomials}
\input{e_polynomials.tex}

\newpage
\chapter{Main results}
\input{results.tex}



\chapter{Conclusiones}




\addcontentsline{toc}{chapter}{Bibliografía}
\bibliographystyle{plain}
\bibliography{bibliografia}



\newpage
\addcontentsline{toc}{chapter}{Anexos}
\chapter*{Anexos}

Agregar en esta sección los anexos de ser necesario.


\end{document}
